\documentclass[11pt]{article}
\usepackage[margin=1in]{geometry}
\usepackage{lmodern}
\usepackage[T1]{fontenc}
\usepackage{microtype}
\usepackage{amsmath,amssymb,mathtools,bm}
\usepackage{booktabs}
\usepackage{graphicx}
\usepackage{caption}
\usepackage{float}
\usepackage{enumitem}
\usepackage{xcolor}
\usepackage[colorlinks=true,linkcolor=blue!50!black,citecolor=blue!50!black,urlcolor=blue!50!black]{hyperref}

\captionsetup{font=small,labelfont=bf}
\setlist{itemsep=2pt,topsep=4pt}

\newcommand{\E}{\mathbb{E}}
\newcommand{\Var}{\operatorname{Var}}
\newcommand{\Cov}{\operatorname{Cov}}
\newcommand{\bx}{\bm{x}}
\newcommand{\bb}{\bm{\beta}}
\newcommand{\dB}{\,\textrm{dB}}

\title{From fitted Hawkes parameters to interpretable quantities:\\
noise, missing calls and the drivers of the call rate}
\author{Rob Schick\\[2pt]\small Draft for discussion, not for circulation}
\date{28 October 2026}

\begin{document}
\maketitle

\begin{abstract}
\noindent
We fit a log-Gaussian Cox background with self-exciting (Hawkes) triggering to
North Atlantic right whale upcall detections at three passive acoustic buoys
in southern New England. This note describes three derived quantities we
propose to report, each a functional of the posterior. They are the
multiplicative effect of a fixed increase in ambient noise on the detected
call rate; a counterfactual count of calls ``missing'' at noisy times; and a
day-by-day multiplicative decomposition of the background rate. We give the
definitions, the computations, the current estimates and the assumptions each
one leans on. We also set out three extensions we would like feedback on: a
sonar-equation benchmark for the noise slope, a noise-dependent branching
ratio as a model-based handle on masking versus behavior, and an honest
attribution of the seasonal trend.
\end{abstract}

\paragraph{Status of the fits.}
COX01 is the clean-data refit (100k iterations, burn-in 75k). NS01 is the
chain extended to 150k with burn-in 100k; its $-2\log L$ trace is flat over
the retained draws. NS02 is still the 100k chain with burn-in 50k. That chain
had not converged, and a 150k extension is running, so the NS02 numbers below
are provisional. Every functional is computed from 2{,}000 equally spaced
post-burn-in draws.

\section{Model}\label{sec:model}

Let $t_1<\dots<t_n$ be detected upcall times (minutes) on $[0,T]$, with $T$
covering 1 October 2021 to 1 March 2022 (151 days). We observe $n=6{,}098$
calls at COX01, $13{,}017$ at NS01 and $33{,}611$ at NS02. The conditional
intensity is
\begin{equation}\label{eq:lambda}
\lambda(t\mid\mathcal{H}_t)=\mu(t)+\sum_{t_i<t}\alpha\,e^{-\eta(t-t_i)},
\qquad
\mu(t)=\exp\{\bx(t)^\top\bb+\delta\,W(t)\}.
\end{equation}
The kernel is unnormalized, so the branching ratio (the expected number of
calls directly triggered by one call) is $m=\alpha/\eta$. We need $m<1$. The
expected cluster size is $1/(1-m)$, and the half-life of the excitation is
$\log 2/\eta$.

\paragraph{Background.}
$\bx(t)$ contains an intercept; ambient noise $x_n(t)$ (Section~\ref{sec:db});
daily sea-surface temperature $x_{s}(t)$ (NOAA OISST, standardized); and sine
and cosine pairs at periods of 1, 7, 14, 30 and 60 days. That makes
$p=13$ columns. $W$ is a zero-mean, unit-variance Gaussian process with
exponential covariance $\Cov\{W(s),W(s')\}=\exp(-|s-s'|/\rho)$, an
Ornstein--Uhlenbeck process. The range is fixed at $\rho=60$ min, so the
correlation is $e^{-3}$ at 3 h. $\delta>0$ scales it. The background is
evaluated at knots $s_0,\dots,s_M$ spaced 20 min apart
($M\approx10{,}800$) and interpolated linearly between them. The compensator
$\int_0^T\mu$ is computed by the trapezoid rule on the same knots.

\paragraph{Priors and computation.}
$\bb\sim N(\bm 0,100\,I)$, $\log\delta\sim N(0,100)$,
$\alpha\sim\mathrm{Gamma}(0.001,0.001)$ and
$\eta\sim U(0.15,\,3/\min_i(t_{i+1}-t_i))$ per minute. The sampler augments
the data with the latent branching structure (each call is either an
immigrant or the offspring of a specific earlier call), drawn by Gibbs. Given
the branching structure, $\alpha$ is conjugate, with shape $0.001+n_{\rm off}$
and rate $0.001+\eta^{-1}\sum_i(1-e^{-\eta(T-t_i)})$. $\bb$, $\log\delta$ and
$\eta$ are updated by adaptive random-walk Metropolis (target acceptance
0.234), and $W$ at the knots by elliptical slice sampling.

\paragraph{A posterior predictive check.}
Under~\eqref{eq:lambda}, ignoring edge effects (the half-life is about
2~min, far shorter than any interval we sum over), the expected number of
calls in a set $A$ is
\begin{equation}\label{eq:EN}
\E N(A)\approx\frac{1}{1-m}\int_A\mu(t)\,dt.
\end{equation}
Over the whole window the posterior median of~\eqref{eq:EN} is $6{,}086$,
$13{,}014$ and $33{,}589$, against observed counts of $6{,}098$, $13{,}017$
and $33{,}611$. That agreement is a property of a well-mixed fit, not
evidence for the noise model, but it confirms the functionals below are on
the right scale. Table~\ref{tab:params} gives the parameter summaries.

\begin{table}[H]
\centering\small
\caption{Posterior medians (95\% equal-tailed intervals) of the parameters
used below. $\beta_n$ and $\beta_s$ are coefficients on the standardized
noise and SST covariates. NS02 is provisional.}
\label{tab:params}
\begin{tabular}{lccc}
\toprule
& COX01 & NS01 & NS02 \\
\midrule
$\alpha$ & 0.203 (0.185, 0.223) & 0.223 (0.209, 0.238) & 0.294 (0.283, 0.305)\\
$\eta$ (min$^{-1}$) & 0.340 (0.308, 0.379) & 0.386 (0.361, 0.419) & 0.494 (0.453, 0.527)\\
$m=\alpha/\eta$ & 0.594 (0.561, 0.634) & 0.576 (0.546, 0.604) & 0.596 (0.567, 0.641)\\
$1/(1-m)$ & 2.47 (2.28, 2.73) & 2.36 (2.20, 2.53) & 2.47 (2.31, 2.79)\\
$\log2/\eta$ (min) & 2.04 (1.83, 2.25) & 1.79 (1.66, 1.92) & 1.40 (1.32, 1.53)\\
$\delta$ & 2.03 (1.89, 2.15) & 1.81 (1.73, 1.90) & 1.60 (1.54, 1.66)\\
$\beta_n$ & $-0.590$ ($-0.705$, $-0.466$) & $-1.477$ ($-1.628$, $-1.333$) & $-0.796$ ($-0.889$, $-0.727$)\\
$\beta_s$ & $-2.84$ ($-3.12$, $-2.54$) & $-1.76$ ($-1.97$, $-1.62$) & $-1.39$ ($-1.49$, $-1.31$)\\
\bottomrule
\end{tabular}
\end{table}

\section{Putting the noise coefficient on the decibel scale}\label{sec:db}

The noise covariate is the RMS noise level from the recorder, rescaled per
deployment and then standardized over the knots:
$x_n(s)=\{y(s)-\bar y\}/\mathrm{sd}(y)$. The RMS level in
dB re 1\,$\mu$Pa is an affine function of $y$; a regression on the raw
one-minute files recovers it with $R^2>0.999$ at every buoy. So
\[
x_n(s)=\frac{L(s)-\bar L}{s_L},
\]
where $L$ is the noise level in dB and $s_L$ is its standard deviation over
the knots: 4.00\dB\ (COX01), 6.02\dB\ (NS01) and 5.30\dB\ (NS02). The
coefficient $\beta_n/s_L$ is therefore the change in $\log\mu$ per dB of
noise. The RMS level is measured in the right whale upcall band, at the
hydrophone.

\section{Fig.~4: the effect of a fixed increase in noise}\label{sec:fig4}

\paragraph{Estimand.}
For an increase of $\Delta$ dB at a fixed time, holding everything else fixed,
\begin{equation}\label{eq:R}
R(\Delta)=\frac{\mu(t;\,L+\Delta)}{\mu(t;\,L)}=\exp\!\left(\frac{\beta_n\Delta}{s_L}\right),
\end{equation}
which doesn't depend on $t$. By~\eqref{eq:EN}, $R(\Delta)$ is also the ratio of
expected \emph{total} detected calls, triggered calls included, as long as noise
acts only on $\mu$. Section~\ref{sec:thin} questions that assumption. We take
$\Delta=5$\dB, which is about one $s_L$ at each buoy, so the comparison is
between sites on the same physical scale rather than the same standardized
one. The posterior of $R$ comes from the draws of $\beta_n$, with $s_L$
treated as known.

\begin{table}[H]
\centering\small
\caption{Noise effect on the detected call rate. The slope is $\beta_n/s_L$,
in log-rate per dB. $R(\Delta)$ is the multiplicative effect of $+\Delta$\dB.}
\label{tab:noise}
\begin{tabular}{lcccc}
\toprule
& $s_L$ (dB) & slope (per dB) & $R(1\dB)$ & $R(5\dB)$\\
\midrule
COX01 & 4.00 & $-0.148$ ($-0.176$, $-0.117$) & 0.863 (0.839, 0.890) & 0.478 (0.415, 0.558)\\
NS01 & 6.02 & $-0.245$ ($-0.271$, $-0.221$) & 0.782 (0.763, 0.801) & 0.293 (0.259, 0.331)\\
NS02 & 5.30 & $-0.150$ ($-0.168$, $-0.137$) & 0.861 (0.846, 0.872) & 0.472 (0.433, 0.504)\\
\bottomrule
\end{tabular}
\end{table}

\begin{figure}[H]
\centering
\includegraphics[width=\textwidth]{volume_knob.pdf}
\caption{Fig.~4 (talk version). Each row shows $10\,R(5\dB)$ rounded to whole
calls. The text gives the posterior median and 95\% interval of $R(5\dB)$.}
\end{figure}

\paragraph{Proposal: a sonar-equation benchmark for the slope.}
A per-dB slope can be compared with what pure masking predicts. In the passive
sonar equation, a call is detected when $\mathrm{SL}-\mathrm{TL}(r)-L\ge
\mathrm{DT}$, with transmission loss $\mathrm{TL}(r)=k\log_{10}r$. The maximum
detection range is then $r_{\max}=10^{(\mathrm{SL}-L-\mathrm{DT})/k}$. If the
detected call rate is proportional to the area within $r_{\max}$,
\begin{equation}\label{eq:bench}
\frac{\partial\log\mu}{\partial L}=-\frac{2\ln 10}{k}
=\begin{cases}-0.230 & k=20\ \text{(spherical)},\\ -0.307 & k=15\ \text{(``practical'')},\\ -0.461 & k=10\ \text{(cylindrical)}.\end{cases}
\end{equation}

\paragraph{Assumptions behind~\eqref{eq:bench}.}
\begin{enumerate}[label=A\arabic*.,leftmargin=2.2em]
\item \emph{Detection depends only on the signal-to-noise ratio at the
  hydrophone.} Masking happens at the receiver, and $L$ is measured at the
  receiver in the upcall band, so the covariate is the right quantity. This
  is the strongest assumption in the argument.
\item \emph{Transmission loss is a power law}, $k\log_{10}r$, the same in all
  directions, with negligible absorption at upcall frequencies. Shallow-water
  loss usually lies between cylindrical and spherical, and $k$ varies with site
  and season.
\item \emph{Source level does not depend on noise.} SL and DT may vary from
  call to call with any distribution, provided it is independent of $L$:
  every call's $r_{\max}$ then scales by the same factor $10^{-\Delta/k}$.
\item \emph{Callers are spread uniformly in the plane, out beyond $r_{\max}$.}
\item \emph{The calling rate per animal does not depend on noise.}
\end{enumerate}

\paragraph{The shape of the detection function does not matter.}
Let detection be probabilistic, with probability $g(\mathrm{SNR})$ for any
function $g$. Under A2--A4 the detected area is
\[
\int_0^\infty g\bigl(\mathrm{SL}-\mathrm{DT}-L-k\log_{10}r\bigr)\,2\pi r\,dr
=10^{2(\mathrm{SL}-\mathrm{DT}-L)/k}\int_0^\infty g\bigl(-k\log_{10}\rho\bigr)\,2\pi\rho\,d\rho,
\]
by the substitution $r=10^{(\mathrm{SL}-\mathrm{DT}-L)/k}\rho$. The remaining
integral doesn't involve $L$, so the slope is exactly $-2\ln10/k$ whatever the
detection curve. A smooth detection curve only matters in combination with a
failure of A2 or A4.

\paragraph{Relaxing the assumptions.}
Three departures each change the slope in a known way.
\begin{itemize}
\item \emph{Lombard compensation} (fails A3): if callers raise their source
  level by $c$ dB per dB of noise, $r_{\max}\propto10^{-(1-c)L/k}$. Tagged right
  whales show $c\approx1$ at moderate noise (Parks et al.\ 2011).
\item \emph{A density gradient} (fails A4): if density falls with distance as
  $r^{-a}$, $0\le a<2$, the detected rate scales as $r_{\max}^{2-a}$. If
  $r_{\max}$ exceeds the area callers occupy, the slope tends to 0 at low
  noise (saturation).
\item \emph{Calling less in noise} (fails A5): if each animal's calling rate is
  proportional to $e^{bL}$, $b$ adds directly to the slope.
\end{itemize}
Together,
\begin{equation}\label{eq:slope4}
\frac{\partial\log\mu}{\partial L}\approx-\frac{(1-c)(2-a)\ln10}{k}+b,
\end{equation}
one equation in four unknowns $(k,a,c,b)$.

\paragraph{What the slopes say.}
\begin{itemize}
\item NS01's slope, $-0.245$ ($-0.271$, $-0.221$), contains the spherical
  benchmark. It is consistent with pure masking at $k=20$ ($c=a=b=0$). It is
  equally consistent with some compensation offset by some reduction in calling.
\item COX01 and NS02 are shallower ($\approx-0.15$). Under
  \eqref{eq:slope4} there are three readings:
  \begin{itemize}
  \item partial compensation, $c\approx0.35$ at $k=20$ or $c\approx0.5$ at $k=15$;
  \item callers concentrated near the buoy, $a>0$;
  \item steeper effective transmission loss: pure masking with $k\approx31$
    gives $-0.15$.
  \end{itemize}
\item A slope \emph{steeper} than any plausible masking slope would require
  $b<0$: fewer calls made, not just fewer heard. No site shows that for
  $k\le20$.
\end{itemize}

\paragraph{From a frame to a test.}
We don't offer~\eqref{eq:bench} as a test yet. Three additions would make it
one:
\begin{itemize}
\item site-specific $k$ from propagation modeling;
\item a low-dimensional spline in $L$ in place of the linear term, to detect
  saturation (curvature at low noise) directly;
\item some handle on the density gradient $a$, for example from array
  localizations.
\end{itemize}
With $k$ fixed, \eqref{eq:slope4} bounds the combination of $c$ and $b$. Two
effects attenuate the fitted slope toward 0 and should be kept in mind. The
covariate is a 1-minute RMS averaged to 20-minute knots, while masking acts
call by call, so short noise spikes behave like measurement error. Orphaned
answers (Section~\ref{sec:thin}) inflate the apparent background most at
noisy times.

\section{Fig.~6: calls missing at noisy times}\label{sec:fig6}

\paragraph{Estimand.}
Let $x_q$ be the $q$-quantile of $x_n$ over the knots of that buoy ($q=0.1$).
The counterfactual background caps the noise at $x_q$ and leaves quieter
times unchanged:
\begin{equation}\label{eq:cap}
\mu_q(s)=\mu(s)\exp\bigl[\beta_n\{\min(x_n(s),x_q)-x_n(s)\}\bigr]\ \ge\ \mu(s).
\end{equation}
For day $d$ let $\Lambda_d=\int_d\mu$ and $\Lambda^q_d=\int_d\mu_q$ (trapezoid
rule on the knots). Using~\eqref{eq:EN}, the expected detected calls are
$\Lambda_d/(1-m)$ as fitted and $\Lambda^q_d/(1-m)$ in the quiet
counterfactual. The gap splits into missing immigrant (``contact'') calls,
$\Lambda^q_d-\Lambda_d$, and the calls those would have triggered,
$(\Lambda^q_d-\Lambda_d)\,m/(1-m)$. Every quantity is computed per posterior
draw, jointly in $(\bb,\delta,W,\alpha,\eta)$, and then summarized. These are
posteriors of \emph{expected} counts, not predictive distributions, so no
Poisson or cluster variability is added.

\begin{table}[H]
\centering\small
\caption{Season totals, posterior medians (95\% intervals). ``Missing'' is
the quiet total minus the fitted total. The share is the ratio of the two
posterior medians. Noise levels are on the knots.}
\label{tab:quiet}
\footnotesize\setlength{\tabcolsep}{4pt}
\begin{tabular}{lccccc}
\toprule
& \shortstack{$L_{0.1}$ / median $L$\\(dB)} & fitted & quiet ($q=0.1$) & missing & share\\
\midrule
COX01 & 72.9 / 77.6 & 6{,}086 (5{,}737, 6{,}466) & 11{,}456 (9{,}830, 13{,}119) & 5{,}357 (3{,}856, 6{,}933) & 47\%\\
NS01 & 77.6 / 84.2 & 13{,}014 (12{,}477, 13{,}541) & 30{,}738 (27{,}494, 34{,}533) & 17{,}733 (14{,}698, 21{,}304) & 58\%\\
NS02 & 76.8 / 82.5 & 33{,}589 (32{,}749, 34{,}552) & 87{,}399 (79{,}577, 98{,}646) & 53{,}667 (46{,}244, 64{,}916) & 61\%\\
\bottomrule
\end{tabular}
\end{table}

\begin{figure}[H]
\centering
\includegraphics[width=\textwidth]{quiet_season_ns01.pdf}
\caption{Fig.~6(i). NS01 daily posterior-median expected calls, fitted (gray)
and under the capped-noise counterfactual~\eqref{eq:cap}. The dark band is
missing immigrant calls; the light band is the calls they would have
triggered. Lower panel: season totals for all three buoys.}
\end{figure}

\begin{figure}[H]
\centering
\includegraphics[width=\textwidth]{quiet_season_all.pdf}
\caption{Fig.~6(ii). The same for the three buoys on a shared axis.}
\end{figure}

\paragraph{What the counterfactual does and doesn't assume.}
\begin{itemize}
\item There's no extrapolation in $x_n$: $x_q$ is inside the observed range,
  and the cap lowers the noise at 90\% of the knots.
\item It does lean on the log-linear form across the whole range. If the
  effect saturates at low noise (Section~\ref{sec:fig4}, A4), the
  linear model overstates $\Lambda^q$.
\item Everything else, $W$ included, is held at its fitted value. If $W$ has
  absorbed residual structure correlated with noise, the attribution to noise
  is affected in a direction we can't sign a priori.
\item The gap is ``missed by the detector or never made''. Detections alone
  don't separate masking from a behavioral reduction in calling.
  Section~\ref{sec:thin} suggests one model-based handle.
\end{itemize}

\paragraph{Proposed sensitivity analyses.}
Report the missing share as a function of $q\in\{0.1,0.25,0.5\}$, using the
export's existing \texttt{PARAM\_QUIET\_Q} switch. Repeat the calculation
under a spline-in-$L$ noise effect. As an alternative estimand that doesn't
depend on choosing a reference level, report the average marginal loss
$-\int\partial\E N/\partial L\,dt$ in calls per dB.

\subsection{Thinning, and a handle on masking versus behavior}\label{sec:thin}

If masking is the mechanism, it removes \emph{every} call with some
probability $p(t)$, triggered calls included. Let the true calls follow a
Hawkes process with background $\mu^\ast(t)$ and branching ratio
$m^\ast=\alpha^\ast/\eta$. The detector then sees a thinned Hawkes process.

\paragraph{Assumptions.}
\begin{enumerate}[label=B\arabic*.,leftmargin=2.2em]
\item \emph{Detection is independent across calls, given the noise.} This is
  the critical assumption. If an answer comes from a different whale at an
  unrelated range, it is detected with probability $p(t)$ independently of its
  parent. If it is the same whale repeating in a bout, or a nearby whale, then
  a detected parent makes a detected answer very likely: same source level,
  same range, same noise. The model can't distinguish repeats from replies, so
  we don't know which case dominates.
\item \emph{$p$ is roughly constant within a cluster.} The half-life is about
  2~min, and the noise changes more slowly than that.
\item \emph{The thinned process is approximately Hawkes.} It isn't exactly: a
  detected answer whose parent was missed enters as an apparent immigrant (an
  ``orphan'').
\end{enumerate}

\paragraph{Consequences under B1.}
\begin{enumerate}[label=(\alph*)]
\item A detected call's detected offspring number $p\,m^\ast$ on average, so
  the detected branching ratio falls with noise: $m(t)\approx p(t)\,m^\ast$.
  The current model holds $m$ fixed, so under masking it is misspecified.
\item Relative to the true immigrants, orphans number about $(1-p)\,m^\ast$.
  They inflate the apparent background most at noisy times, which pulls the
  fitted $\beta_n$ toward 0.
\item Both effects make the missing counts in Table~\ref{tab:quiet}
  underestimates. A larger $m$ at quiet times raises the quiet multiplier
  $1/(1-m_q)$, and the true background slope is steeper than fitted.
  Saturation (A4) pushes the other way, so the net direction is not
  guaranteed.
\end{enumerate}

\paragraph{Proposed model.}
Let the noise at the parent's time scale its offspring:
$\alpha_i=\alpha\exp\{\gamma\,x_n(t_i)\}$. Given the branching structure,
call $i$'s offspring count is Poisson with mean
$\alpha e^{\gamma x_i}(1-e^{-\eta(T-t_i)})/\eta$. So $\alpha$ stays conjugate,
with full conditional
\[
\alpha\mid\cdot\sim\mathrm{Gamma}\Bigl(0.001+n_{\rm off},\;
0.001+\eta^{-1}\textstyle\sum_i e^{\gamma x_i}\bigl(1-e^{-\eta(T-t_i)}\bigr)\Bigr),
\]
and the only new piece is a Metropolis step for $\gamma$. If the noise at the
child's time matters more, a variant puts $e^{\gamma x(t)}$ inside the kernel
integral. That is nearly identical here, because $x_n$ barely changes within a
2-minute cluster.

\paragraph{Interpreting $\gamma$.}
If masking is the mechanism:
\begin{itemize}
\item $\gamma\approx\beta_n$: answers are thinned independently of their
  parents (different whales at unrelated ranges);
\item $\gamma\approx0$: answers are detected together with their parents
  (same-whale bouts);
\item $\gamma/\beta_n\in[0,1]$ measures roughly the fraction of answers whose
  detection is independent of their parent's.
\end{itemize}
If noise acts on behavior instead, $\gamma<0$ would mean whales answer less in
noise. So $\gamma$ constrains the mechanism; it doesn't prove masking. Few calls
are detected at high noise, so the posterior for $\gamma$ may be wide exactly
where it matters; its width will tell us.

\section{Fig.~7: what moved the call rate}\label{sec:fig7}

\paragraph{Decomposition.}
At each knot,
\[
\log\mu(s)=\beta_0+\underbrace{\beta_n x_n(s)}_{c_n(s)}+\underbrace{\beta_s x_s(s)}_{c_s(s)}
+\underbrace{\textstyle\sum_j\beta_j h_j(s)}_{c_h(s)}+\underbrace{\delta W(s)}_{c_W(s)}.
\]
For component $c$ and day $d$ let $\bar c(d)$ be the mean over the knots in
$d$. Fig.~7 plots
\[
F_c(d)=\exp\Bigl\{\bar c(d)-\tfrac{1}{D}\textstyle\sum_{d'}\bar c(d')\Bigr\},
\]
the factor by which component $c$ multiplied the background on day $d$,
relative to its own season average. The product over components of $F_c(d)$
times a constant is the geometric-mean background on day $d$. The
components use per-knot posterior medians, a plug-in summary; the median of
the sum is not the sum of the medians, so the panels don't multiply exactly
to the posterior-median background.

\paragraph{Answering.}
The bottom panel is the data-conditional ratio
\[
A(d)=\frac{\sum_{s\in d}\{\mu(s)+\phi(s)\}}{\sum_{s\in d}\mu(s)},\qquad
\phi(s)=\sum_{t_i<s}\alpha e^{-\eta(s-t_i)},
\]
on a 5-min grid, using the observed calls and plug-in posterior medians. It
is left uncentered, so $A=1$ on call-free days. Its geometric mean over days
is 1.48 (COX01), 1.57 (NS01) and 2.11 (NS02). For NS01, $\mu$, $\alpha$ and
$\eta$ in this panel still come from the 100k fit; that will be updated.

\begin{figure}[H]
\centering
\includegraphics[width=\textwidth]{rate_drivers_ns01.pdf}
\caption{Fig.~7, NS01. Daily multiplicative factors $F_c(d)$ for the four
background components, and the answering ratio $A(d)$. Log scale.
COX01 and NS02 are in the Appendix.}
\end{figure}

\paragraph{Variance attribution.}
To summarize the panels we split the day-to-day variance of the daily-mean
log background, $S(d)=\sum_c\bar c(d)$, into
$\Cov\{\bar c(d),S(d)\}/\Var\{S(d)\}$. These shares sum to one and can be
negative when components are negatively correlated
(Table~\ref{tab:var}).

\begin{table}[H]
\centering\small
\caption{Covariance shares of the variance of the log background, at daily
means and at the 20-min knots. Plug-in posterior medians.}
\label{tab:var}
\begin{tabular}{lcccccccc}
\toprule
& \multicolumn{4}{c}{daily means} & \multicolumn{4}{c}{knots}\\
\cmidrule(lr){2-5}\cmidrule(lr){6-9}
& noise & SST & harmonics & GP & noise & SST & harmonics & GP\\
\midrule
COX01 ($\Var S=9.84$) & $-0.02$ & 0.79 & 0.18 & 0.06 & 0.00 & 0.62 & 0.15 & 0.22\\
NS01 ($\Var S=6.90$) & 0.29 & 0.50 & 0.08 & 0.14 & 0.29 & 0.37 & 0.06 & 0.29\\
NS02 ($\Var S=2.71$) & $-0.07$ & 0.49 & 0.14 & 0.44 & 0.02 & 0.28 & 0.10 & 0.60\\
\bottomrule
\end{tabular}
\end{table}

What Table~\ref{tab:var} shows:
\begin{itemize}
\item Noise is a major driver of day-to-day variation only at NS01. A
  component's variance is $\beta^2$ times the variance of its covariate. NS01
  has the largest per-SD effect ($-1.48$) and the most variable daily noise:
  the SD of daily-mean $L$ is 4.9\dB, against 2.0\dB\ at COX01 and 3.7\dB\ at
  NS02.
\item At COX01 and NS02 the noise contribution is negatively correlated with
  the rest of the background ($-0.33$ and $-0.48$ at daily means; $-0.15$ and
  $-0.30$ at knots). That cancels most of its own variance, so its share is
  near zero at both scales, not just day to day.
\item At NS02, daily noise and SST contributions correlate at $-0.74$. Both
  coefficients are negative, so louder days fall later in the season, as
  winter storms build while whales arrive. Much of the variation that
  identifies $\beta_n$ there is sub-seasonal, and the split between
  noise and season is weakly identified.
\item The GP share at NS02 (0.44 daily, 0.60 at knots) says there is a lot of
  residual structure on scales the covariates don't reach. We read it as a
  goodness-of-fit diagnostic as much as a driver.
\end{itemize}

\paragraph{Identifiability of the seasonal trend.}
The harmonics have periods of at most 60 days over a 151-day window. They can
represent recurring rhythms but approximate a monotone trend only poorly. The
GP has a 3-hour memory and can't carry one either. SST is the only smooth,
monotone term, and over the window it correlates with time at $-0.99$,
$-0.98$ and $-0.96$. In practice $c_s$ is therefore the seasonal trend, and
$\beta_s$ should not be read as an effect per $^\circ$C. With one season per
site, temperature and time of year are confounded. We propose either
replacing SST with an explicit smooth trend in time, or reporting $c_s+c_h$
together as ``season'', and keeping SST for when multiple seasons or
decorrelated sites are available.

\paragraph{Proposed improvements.}
Compute $F_c(d)$ and $A(d)$ per draw and show pointwise 95\% bands. Report
the shares in Table~\ref{tab:var} with posterior intervals. Given the
correlations between components, consider a Shapley-type attribution, which
is order-invariant, as an alternative to covariance shares.

\section{Questions for discussion}

\begin{enumerate}
\item Is the sonar-equation benchmark~\eqref{eq:bench} a useful frame for the
  noise slope? How should uncertainty in $k$ and the density gradient $a$ be
  carried into~\eqref{eq:slope4}?
\item Detection belongs in the model. Is a noise-dependent kernel
  (Section~\ref{sec:thin}) adequate, or should we fit the thinning explicitly,
  with undetected calls as latent variables?
\item Is the capped-noise expected count the right estimand for ``calls
  lost to noise''? Or is the average marginal effect in calls per dB cleaner,
  since it avoids choosing a reference level?
\item NS01's $\eta$ drifted across the first 100k iterations, with a matching
  shift between the GP background and the excitation. That is the usual
  clustering-versus-contagion trade-off. Would a reparameterization in
  $(m,\eta)$, or a prior on $\delta$, reduce it?
\end{enumerate}

\appendix
\section{Fig.~7 for COX01 and NS02}

\begin{figure}[H]
\centering
\includegraphics[width=\textwidth]{rate_drivers_cox01.pdf}
\caption{COX01.}
\end{figure}

\begin{figure}[H]
\centering
\includegraphics[width=\textwidth]{rate_drivers_ns02.pdf}
\caption{NS02 (provisional fit).}
\end{figure}

\raggedright
\paragraph{Reference.}
Parks, S.~E., Johnson, M., Nowacek, D.\ and Tyack, P.~L. (2011). Individual
right whales call louder in increased environmental noise.
\emph{Biology Letters} 7(1), 33--35. doi:10.1098/rsbl.2010.0451.

\paragraph{Reproducibility.}
Cluster export: \texttt{10b\_paramExport.R}. Figures:
\texttt{14\_volumeKnob.R}, \texttt{15\_quietSeason.R} and
\texttt{17\_rateDrivers.R}, with the shared helpers in
\texttt{src/talk\_figs.R}.

\end{document}
